\chapter{Limieten en continuïteit}\label{ch:g12:limcont}

Dit hoofdstuk breidt het begrip limiet uit van \omterm{def:g12:seq:sequence}{rijen} naar \omterm{def:g11:func:function}{functies} van een
reële veranderlijke, en voert de \omterm{def:g12:limcont:continuity}{continuïteit} in. Het centrale resultaat is de
tussenwaardestelling, die kwalitatieve informatie (``$f$ is \omterm{def:g12:limcont:continuity}{continu} en
verandert van teken'') omzet in het bestaan van oplossingen van
\omterm{def:g10:algebra:equation}{vergelijkingen}.

\section{Limieten van functies}

Overal in dit hoofdstuk is $f$ een \omterm{def:g11:func:function}{functie} gedefinieerd op een \omterm{def:g10:numbers:interval}{interval} of een
\omterm{def:g10:numbers:interunion}{vereniging} van intervallen $D \subseteq \R$.

\begin{definition}[Limiet in oneindig]\label{def:g12:limcont:liminf}
Zij $f$ gedefinieerd op een \omterm{def:g10:numbers:interval}{interval} $\intco{a}{+\infty}$ en $\ell \in \R$.
We zeggen dat $f$ \emph{naar $\ell$ streeft in $+\infty$}\index{limiet!van een
functie}, geschreven $\lim\limits_{x\to+\infty} f(x) = \ell$, wanneer elk
open \omterm{def:g10:numbers:interval}{interval} dat $\ell$ bevat, alle waarden $f(x)$ bevat voor $x$ groot
genoeg: voor elke $\varepsilon > 0$ bestaat er een $A$ zodat
$\abs{f(x) - \ell} \leq \varepsilon$ voor alle $x \geq A$.

We zeggen dat $f$ naar $+\infty$ streeft in $+\infty$ wanneer er voor elke
$M \in \R$ een $A$ bestaat zodat $f(x) \geq M$ voor alle $x \geq A$. Limieten
in $-\infty$ en limieten gelijk aan $-\infty$ worden analoog gedefinieerd.
\end{definition}

\begin{definition}[Limiet in een punt]\label{def:g12:limcont:limpoint}
Zij $a \in \R$ en zij $f$ gedefinieerd in de buurt van $a$ (behalve
eventueel in $a$ zelf). We zeggen dat $f$ in $a$ naar $\ell \in \R$ streeft,
geschreven $\lim\limits_{x \to a} f(x) = \ell$, wanneer er voor elke
$\varepsilon > 0$ een $\delta > 0$ bestaat zodat
$\abs{f(x) - \ell} \leq \varepsilon$ voor alle $x \in D$ met
$\abs{x - a} \leq \delta$.

We zeggen dat $f$ in $a$ naar $+\infty$ streeft wanneer er voor elke $M$ een
$\delta > 0$ bestaat zodat $f(x) \geq M$ zodra $x \in D$ en
$\abs{x - a} \leq \delta$. Eenzijdige limieten ($x \to a^+$, $x \to a^-$)
beperken $x$ tot $x > a$ of $x < a$.
\end{definition}

\begin{example}
$\lim\limits_{x\to+\infty} \frac{1}{x} = 0$,
$\lim\limits_{x\to 0^+} \frac{1}{x} = +\infty$,
$\lim\limits_{x\to 0^-} \frac{1}{x} = -\infty$. De \omterm{def:g11:func:function}{functie}
$x \mapsto \frac1x$ heeft in $0$ geen (tweezijdige) limiet.
\end{example}

\begin{definition}[Asymptoten]\label{def:g12:limcont:asymptote}
De rechte $y = \ell$ is een \emph{horizontale asymptoot}\index{asymptoot} van
de kromme van $f$ in $+\infty$ (respectievelijk $-\infty$) wanneer
$\lim_{x\to+\infty} f(x) = \ell$ (respectievelijk in $-\infty$). De rechte
$x = a$ is een \emph{verticale asymptoot}\index{asymptoot} wanneer $f$ in $a$
een oneindige eenzijdige limiet heeft.
\end{definition}

\begin{omfigure}
\begin{tikzpicture}
\begin{axis}[omaxis,
  width=8.8cm, height=6cm,
  xmin=-2.9, xmax=4.9, ymin=-5.2, ymax=9.2,
  xtick={1}, ytick={2}]
  \draw[omProp, dashed] (axis cs:-2.9,2) -- (axis cs:4.9,2);
  \draw[omProp, dashed] (axis cs:1,-5.2) -- (axis cs:1,9.2);
  \addplot[omDef, thick, domain=-2.8:0.855, samples=60] {2 + 1/(x-1)};
  \addplot[omDef, thick, domain=1.145:4.8, samples=60] {2 + 1/(x-1)};
\end{axis}
\end{tikzpicture}

{\small De kromme van $f(x) = 2 + \frac{1}{x-1}$ met haar \omterm{def:g12:limcont:asymptote}{horizontale
asymptoot} $y = 2$ (in $\pm\infty$) en haar \omterm{def:g12:limcont:asymptote}{verticale asymptoot} $x = 1$.}
\end{omfigure}

\begin{proposition}[Bewerkingen met limieten]\label{prop:g12:limcont:operations}
De regels van \cref{prop:g12:seq:operations} (sommen, producten, quotiënten)
gelden woordelijk voor limieten van \omterm{def:g11:func:function}{functies}, in een punt of in oneindig, met
dezelfde vier onbepaalde vormen $(+\infty)+(-\infty)$, $0\times\infty$,
$\frac{\infty}{\infty}$ en $\frac{0}{0}$.
\end{proposition}

\begin{proof}
De bewijzen zijn identiek aan het geval van de \omterm{def:g12:seq:sequence}{rijen}, met ``voor $n \geq N$''
vervangen door ``voor $x \geq A$'' of ``voor $\abs{x-a} \leq \delta$''.
\end{proof}

\begin{theorem}[Limiet van een samenstelling]\label{thm:g12:limcont:composition}
Zij $f, g$ \omterm{def:g11:func:function}{functies} en $a, b, c$ \omterm{def:g10:numbers:sets}{reële getallen} of $\pm\infty$. Is
$\lim\limits_{x \to a} f(x) = b$ en $\lim\limits_{y \to b} g(y) = c$, dan
geldt
\[
\lim_{x \to a} g\bigl(f(x)\bigr) = c .
\]
Dezelfde uitspraak geldt met een \omterm{def:g12:seq:sequence}{rij} in de plaats van $f$: is $u_n \to b$ en
$\lim_{y\to b} g(y) = c$, dan is $g(u_n) \to c$.
\end{theorem}

\begin{proof}
Neem $b, c$ eindig (de andere gevallen verlopen analoog). Zij
$\varepsilon > 0$. Er bestaat een $\delta > 0$ zodat uit
$\abs{y - b} \leq \delta$ volgt dat $\abs{g(y) - c} \leq \varepsilon$; en er
bestaat een $\delta' > 0$ zodat uit $\abs{x - a} \leq \delta'$ volgt dat
$\abs{f(x) - b} \leq \delta$. De twee aan elkaar rijgen geeft
$\abs{g(f(x)) - c} \leq \varepsilon$ voor $\abs{x - a} \leq \delta'$.
\end{proof}

\begin{theorem}[Vergelijking en insluiting]\label{thm:g12:limcont:squeeze}
Zij $f, g, h$ \omterm{def:g11:func:function}{functies} en $a \in \R \cup \{\pm\infty\}$.
\begin{enumerate}
  \item Is $f \leq g$ in de buurt van $a$ en $f(x) \to +\infty$ voor
        $x \to a$, dan is $g(x) \to +\infty$.
  \item Is $f \leq g \leq h$ in de buurt van $a$ en streven $f$ en $h$ in $a$
        naar dezelfde eindige limiet $\ell$, dan is $g(x) \to \ell$ voor
        $x \to a$.
\end{enumerate}
\end{theorem}

\begin{proof}
Hetzelfde argument als voor \omterm{def:g12:seq:sequence}{rijen} (\cref{thm:g12:seq:squeeze}).
\end{proof}

\begin{method}[Limieten berekenen in de praktijk]\label{met:g12:limcont:limits}
\begin{itemize}
  \item Veeltermen en rationale \omterm{def:g11:func:function}{functies} in $\pm\infty$: haal de hoogste
        macht van $x$ buiten haakjes; de limiet is die van de verhouding van
        de hoogstegraadstermen.
  \item $\frac{0}{0}$ in een punt $a$: haal $(x-a)$ uit teller en noemer, of
        herken een \omterm{def:g11:deriv:derivative}{afgeleide} $\lim_{x\to a}\frac{f(x)-f(a)}{x-a} = f'(a)$
        (\cref{ch:g12:deriv}).
  \item Vierkantswortels: vermenigvuldig met de toegevoegde uitdrukking.
  \item Schommelende factoren ($\cos$, $\sin$): gebruik de insluitstelling.
\end{itemize}
\end{method}

\section{Continuïteit}

\begin{definition}[Continuïteit]\label{def:g12:limcont:continuity}
Zij $f$ gedefinieerd op een \omterm{def:g10:numbers:interval}{interval} $I$ en $a \in I$. De \omterm{def:g11:func:function}{functie} $f$ is
\emph{continu in $a$}\index{continuïteit} wanneer
$\lim\limits_{x \to a} f(x) = f(a)$. Ze is \emph{continu op
$I$}\index{continu} wanneer ze in elk punt van $I$ continu is.
\end{definition}

Intuïtief kun je de \omterm{def:g10:functions:graph}{grafiek} van een continue \omterm{def:g11:func:function}{functie} op een \omterm{def:g10:numbers:interval}{interval} tekenen
zonder de pen op te heffen.

\begin{proposition}[Continuïteit van de gebruikelijke functies]\label{prop:g12:limcont:usual}
Veeltermen, rationale \omterm{def:g11:func:function}{functies}, $x \mapsto \sqrt{x}$, $x \mapsto \abs{x}$,
$\cos$, $\sin$, $\exp$ en $\ln$ zijn \omterm{def:g12:limcont:continuity}{continu} op elk \omterm{def:g10:numbers:interval}{interval} binnen hun
\omterm{def:g11:func:function}{domein}. Sommen, producten, quotiënten (waar ze gedefinieerd zijn) en
samenstellingen van continue \omterm{def:g11:func:function}{functies} zijn \omterm{def:g12:limcont:continuity}{continu}.
\end{proposition}

\begin{proof}[Gedeeltelijk bewijs]
De stabiliteit onder sommen, producten, quotiënten en samenstelling volgt uit
de overeenkomstige stellingen over limieten. De \omterm{def:g12:limcont:continuity}{continuïteit} van
$x \mapsto x$ en van de constante \omterm{def:g11:func:function}{functies} volgt onmiddellijk uit de
definitie, waardoor veeltermen en rationale \omterm{def:g11:func:function}{functies} \omterm{def:g12:limcont:continuity}{continu} zijn. De
\omterm{def:g12:limcont:continuity}{continuïteit} van de overige gebruikelijke \omterm{def:g11:func:function}{functies} wordt bewezen in de
hoofdstukken waar ze opgebouwd worden, of aangenomen.
\end{proof}

\begin{example}
De \emph{entierfunctie} $x \mapsto \floor{x}$ (het grootste gehele getal
$\leq x$) is in geen enkel \omterm{def:g10:numbers:sets}{geheel getal} $n$ \omterm{def:g12:limcont:continuity}{continu}: haar linkerlimiet daar
is $n - 1$ terwijl $\floor{n} = n$.
\end{example}

\begin{omfigure}
\begin{tikzpicture}
\begin{axis}[omaxis,
  width=7.2cm, height=5.6cm,
  xmin=-1.9, xmax=3.5, ymin=-2.5, ymax=2.7,
  xtick={-1,1,2,3}, ytick={-2,-1,1,2}]
  \draw[omDef, thick] (axis cs:-1.5,-2) -- (axis cs:-1,-2);
  \draw[omDef, thick] (axis cs:-1,-1) -- (axis cs:0,-1);
  \draw[omDef, thick] (axis cs:0,0) -- (axis cs:1,0);
  \draw[omDef, thick] (axis cs:1,1) -- (axis cs:2,1);
  \draw[omDef, thick] (axis cs:2,2) -- (axis cs:3,2);
  \fill[omDef] (axis cs:-1,-1) circle (1.5pt);
  \fill[omDef] (axis cs:0,0) circle (1.5pt);
  \fill[omDef] (axis cs:1,1) circle (1.5pt);
  \fill[omDef] (axis cs:2,2) circle (1.5pt);
  \draw[omDef, fill=white] (axis cs:-1,-2) circle (1.5pt);
  \draw[omDef, fill=white] (axis cs:0,-1) circle (1.5pt);
  \draw[omDef, fill=white] (axis cs:1,0) circle (1.5pt);
  \draw[omDef, fill=white] (axis cs:2,1) circle (1.5pt);
  \draw[omDef, fill=white] (axis cs:3,2) circle (1.5pt);
\end{axis}
\end{tikzpicture}

{\small De entierfunctie springt in elk \omterm{def:g10:numbers:sets}{geheel getal}: de waarde (volle stip)
verschilt van de linkerlimiet (open stip).}
\end{omfigure}

\begin{proposition}[Rijen en continuïteit]\label{prop:g12:limcont:seqcont}
Is $u_n \to \ell$ en is $f$ \omterm{def:g12:limcont:continuity}{continu} in $\ell$, dan is $f(u_n) \to f(\ell)$.
In het bijzonder geldt: \omterm{def:g12:seq:limit}{convergeert} een \omterm{def:g12:seq:sequence}{rij} bepaald door $u_{n+1} = f(u_n)$
naar $\ell$ en is $f$ \omterm{def:g12:limcont:continuity}{continu} in $\ell$, dan is $f(\ell) = \ell$.
\end{proposition}

\begin{proof}
De eerste bewering is \cref{thm:g12:limcont:composition}, toegepast met de
\omterm{def:g12:seq:sequence}{rij} $(u_n)$. Voor de tweede ga je in beide leden van $u_{n+1} = f(u_n)$ over
naar de limiet: het linkerlid streeft naar $\ell$, het rechterlid naar
$f(\ell)$, en limieten zijn uniek.
\end{proof}

\section{De tussenwaardestelling}

\begin{theorem}[Tussenwaardestelling]\label{thm:g12:limcont:ivt}
\index{tussenwaardestelling}
Zij $f$ \omterm{def:g12:limcont:continuity}{continu} op $\intcc{a}{b}$ en zij $k$ een willekeurig getal tussen
$f(a)$ en $f(b)$. Dan bestaat er een $c \in \intcc{a}{b}$ met $f(c) = k$.
\end{theorem}

\begin{omfigure}
\begin{tikzpicture}
\begin{axis}[omaxis,
  width=8.8cm, height=6cm,
  xmin=-0.4, xmax=4.4, ymin=-0.5, ymax=4.7,
  xtick={0.5,3.19,3.8}, xticklabels={$a$,$c$,$b$},
  ytick={1.26,2.5,3.98}, yticklabels={$f(a)$,$k$,$f(b)$}]
  \addplot[omProp, dashed, domain=0:3.19] {2.5};
  \addplot[omDef, thick, domain=0.5:3.8] {0.08*x^3 - 0.5*x + 1.5};
  \draw[black, dashed, thin] (axis cs:0.5,0) -- (axis cs:0.5,1.26)
    -- (axis cs:0,1.26);
  \draw[black, dashed, thin] (axis cs:3.8,0) -- (axis cs:3.8,3.98)
    -- (axis cs:0,3.98);
  \draw[black, dashed, thin] (axis cs:3.19,2.5) -- (axis cs:3.19,0);
  \fill (axis cs:0.5,1.26) circle (1.5pt);
  \fill (axis cs:3.8,3.98) circle (1.5pt);
  \fill (axis cs:3.19,2.5) circle (1.5pt);
\end{axis}
\end{tikzpicture}

{\small De tussenwaardestelling: een continue kromme die $(a, f(a))$ met
$(b, f(b))$ verbindt, moet elke horizontale rechte $y = k$ tussen $f(a)$ en
$f(b)$ kruisen.}
\end{omfigure}

\begin{proof}
Door $f$ door $f - k$ te vervangen (en zo nodig door $k - f$) mogen we
aannemen dat $f(a) \leq 0 \leq f(b)$ en zoeken we een nulpunt van $f$. We
bouwen twee \omterm{def:g12:seq:sequence}{rijen} met \emph{bisectie}\index{bisectie}: stel $a_0 = a$,
$b_0 = b$; gegeven $a_n \leq b_n$ met $f(a_n) \leq 0 \leq f(b_n)$, stel
$m = \frac{a_n + b_n}{2}$ en definieer
\[
(a_{n+1}, b_{n+1}) =
\begin{cases}
(a_n, m) & \text{als } f(m) \geq 0,\\
(m, b_n) & \text{als } f(m) < 0.
\end{cases}
\]
In beide gevallen is $f(a_{n+1}) \leq 0 \leq f(b_{n+1})$, is de \omterm{def:g12:seq:sequence}{rij} $(a_n)$
\omterm{def:g12:seq:monotonic}{stijgend}, is $(b_n)$ \omterm{def:g10:functions:variations}{dalend}, en is $b_n - a_n = \frac{b-a}{2^n} \to 0$: de
\omterm{def:g12:seq:sequence}{rijen} zijn aangrenzend (zie \cref{exo:g12:seq:10}) en convergeren naar een
gemeenschappelijke limiet $c \in \intcc{a}{b}$. Door de \omterm{def:g12:limcont:continuity}{continuïteit} is
$f(a_n) \to f(c)$ en $f(b_n) \to f(c)$; overgaan naar de limiet in
$f(a_n) \leq 0$ en $f(b_n) \geq 0$ geeft $f(c) \leq 0 \leq f(c)$, dus
$f(c) = 0$.
\end{proof}

\begin{remark}
Het bewijs met \omterm{thm:g12:limcont:ivt}{bisectie} is een algoritme: het levert willekeurig nauwkeurige
\omterm{def:g10:numbers:approx}{benaderingen} van een oplossing, met bij elke stap een halvering van de fout.
Zo kan een computer $f(x) = 0$ oplossen terwijl hij alleen $f$ kan evalueren.
\end{remark}

\begin{theorem}[Bijectiestelling]\label{thm:g12:limcont:bijection}
Zij $f$ \omterm{def:g12:limcont:continuity}{continu} en \emph{strikt monotoon} op een \omterm{def:g10:numbers:interval}{interval} $I$. Dan heeft de
\omterm{def:g10:algebra:equation}{vergelijking} $f(x) = k$ voor elke $k$ die strikt tussen de limieten van $f$
in de uiteinden van $I$ ligt, een \emph{unieke} oplossing in $I$.
\end{theorem}

\begin{proof}
Bestaan: kies $a, b \in I$ zodat $k$ tussen $f(a)$ en $f(b)$ ligt (dat kan,
want de waarden van $f$ naderen haar limieten in de uiteinden), en pas de
tussenwaardestelling toe op $\intcc{a}{b}$. Uniciteit: is $x < y$, dan geeft
de strikte monotonie dat $f(x) \neq f(y)$, zodat twee verschillende punten
niet allebei oplossing kunnen zijn.
\end{proof}

\begin{method}[$f(x) = k$ oplossen met de bijectiestelling]\label{met:g12:limcont:bijection}
Om te bewijzen dat $f(x) = k$ precies één oplossing heeft in een \omterm{def:g10:numbers:interval}{interval}:
\begin{enumerate}
  \item stel een \omterm{def:g10:functions:table}{verlooptabel} van $f$ op (teken van $f'$, zie
        \cref{ch:g12:deriv});
  \item vergelijk $k$ op elk \omterm{def:g10:numbers:interval}{interval} waar $f$ \omterm{def:g12:limcont:continuity}{continu} en strikt monotoon is
        met de waarden of limieten in de uiteinden;
  \item pas de bijectiestelling op elk zulk \omterm{def:g10:numbers:interval}{interval} toe en tel de
        oplossingen.
\end{enumerate}
Een rekenmachine of \omterm{thm:g12:limcont:ivt}{bisectie} lokaliseert daarna elke oplossing zo fijn als je
wilt.
\end{method}

\begin{example}\label{ex:g12:limcont:cubic}
De \omterm{def:g10:algebra:equation}{vergelijking} $x^3 + x = 1$ heeft precies één reële oplossing. Inderdaad is
$f(x) = x^3 + x$ \omterm{def:g12:limcont:continuity}{continu} en strikt \omterm{def:g12:seq:monotonic}{stijgend} op $\R$ (als som van strikt
\omterm{def:g11:func:monotone}{stijgende functies}), met $\lim_{-\infty} f = -\infty$ en
$\lim_{+\infty} f = +\infty$, zodat de bijectiestelling met $k = 1$ van
toepassing is. Omdat $f(0.6) = 0.816 < 1$ en $f(0.7) = 1.043 > 1$, ligt de
oplossing in $\intoo{0.6}{0.7}$.
\end{example}

\section{Oefeningen}

\begin{exercise}[$\star$]\label{exo:g12:limcont:1}
Bereken de volgende limieten:
\[
\lim_{x\to+\infty} \frac{2x^3 - x + 1}{5x^3 + x^2}, \qquad
\lim_{x\to-\infty} \bigl(x^2 + 3x - 1\bigr), \qquad
\lim_{x\to 2^+} \frac{x+1}{x-2}, \qquad
\lim_{x\to+\infty} \frac{\cos x}{x}.
\]
\end{exercise}

\begin{exercise}[$\star$]\label{exo:g12:limcont:2}
Zij $f(x) = \dfrac{2x^2 - 3x + 1}{x - 1}$, gedefinieerd voor $x \neq 1$.
\begin{enumerate}
  \item Toon aan dat $f(x) = 2x - 1$ voor alle $x \neq 1$.
  \item Leid $\lim\limits_{x \to 1} f(x)$ af. Kan $f$ uitgebreid worden tot
        een \omterm{def:g11:func:function}{functie} die \omterm{def:g12:limcont:continuity}{continu} is in $1$?
\end{enumerate}
\end{exercise}

\begin{exercise}[$\star$]\label{exo:g12:limcont:3}
Bepaal de asymptoten van de kromme van $f(x) = \dfrac{3x - 1}{x + 2}$ en
schets de kromme.
\end{exercise}

\begin{exercise}[$\star\star$]\label{exo:g12:limcont:4}
Bereken
\[
\lim_{x\to+\infty} \bigl(\sqrt{x^2 + x} - x\bigr)
\qquad\text{en}\qquad
\lim_{x\to 0} \frac{\sqrt{1+x} - 1}{x}.
\]
\end{exercise}

\begin{exercise}[$\star\star$]\label{exo:g12:limcont:5}
Zij $f(x) = x + \sin x$.
\begin{enumerate}
  \item Toon aan dat $f(x) \to +\infty$ voor $x \to +\infty$, hoewel $f$ op
        geen enkel \omterm{def:g10:numbers:interval}{interval} $\intco{A}{+\infty}$ strikt monotoon is.
  \item Toon aan dat $g(x) = \dfrac{x + \sin x}{x - \sin x}$ in $+\infty$
        naar $1$ streeft.
\end{enumerate}
\end{exercise}

\begin{exercise}[$\star\star$]\label{exo:g12:limcont:6}
Toon aan dat de \omterm{def:g10:algebra:equation}{vergelijking} $x^3 - 3x + 1 = 0$ precies drie reële
oplossingen heeft, en geef voor elk een \omterm{def:g10:numbers:interval}{interval} van lengte $1$ dat ze
bevat. (Tip: bestudeer het verloop van $f(x) = x^3 - 3x + 1$; haar \omterm{def:g11:deriv:derivative}{afgeleide}
is $3x^2 - 3$.)
\end{exercise}

\begin{exercise}[$\star\star$]\label{exo:g12:limcont:7}
Zij $f \colon \intcc{0}{1} \to \intcc{0}{1}$ \omterm{def:g12:limcont:continuity}{continu}. Toon aan dat $f$ een
\emph{\omterm{pb:g10:functions:1}{vast punt}} heeft: er bestaat een $c \in \intcc{0}{1}$ met $f(c) = c$.
(Tip: pas de tussenwaardestelling toe op $g(x) = f(x) - x$.)
\end{exercise}

\begin{exercise}[$\star\star\star$]\label{exo:g12:limcont:8}
Een wandelaar beklimt een bergpad, vertrekt om 8 uur en komt om 20 uur aan.
De volgende dag daalt ze hetzelfde pad af, opnieuw met vertrek om 8 uur en
aankomst om 20 uur. Toon aan dat er een tijdstip bestaat waarop ze zich op
beide dagen op precies hetzelfde punt van het pad bevindt. (Tip: voer het
\omterm{def:g11:seq:arithmetic}{verschil} van de twee positiefuncties in.)
\end{exercise}

\begin{exercise}[$\star\star\star$]\label{exo:g12:limcont:9}
Zij $f(x) = x^5 + x^3 + x - 1$.
\begin{enumerate}
  \item Toon aan dat $f$ een uniek reëel nulpunt $\alpha$ heeft, en dat
        $\alpha \in \intoo{0}{1}$.
  \item Beschrijf een bisectieprocedure die $\alpha$ op $10^{-6}$ nauwkeurig
        berekent, en geef het aantal iteraties dat daarvoor nodig is,
        vertrekkend van $\intcc{0}{1}$.
\end{enumerate}
\end{exercise}

\section{Opgave: je steekt geen rivier over zonder nat te worden}

\begin{problem}[{Weekendopgave --- de tussenwaardestelling dwingt vaste
punten, antipoden met gelijke temperatuur en stabiele tafels af: drie
onmogelijk klinkende stellingen uit één idee}]
\label{pb:g12:limcont:1}
Verfrommel een kaart van Frankrijk en laat ze ergens op Franse bodem vallen:
een punt van de kaart ligt \emph{precies} boven de plek die het voorstelt. Op
elk ogenblik hebben twee diametraal tegenoverliggende punten van de evenaar
\emph{precies} dezelfde temperatuur. En een wiebelende vierkante tafel op
oneffen grond krijg je altijd stabiel door hem te \emph{draaien}. Drie
salontrucs, één stelling: een continue grootheid kan niet van negatief naar
positief gaan zonder onderweg te verdwijnen (\cref{thm:g12:limcont:ivt}). Deze
opgave oefent limieten, bouwt daarna de drie wonderen --- en weet precies wat
de stelling niet belooft.

\medskip
\noindent\textbf{Deel I --- Vlot met limieten.}
\begin{enumerate}
  \item Bereken $\lim_{x \to +\infty} \dfrac{3x^2 - x}{x^2 + 5}$, \
        $\lim_{x \to +\infty} \dfrac{2x + 1}{x^2 + 1}$ \ en
        $\lim_{x \to +\infty} (x^3 - x^2)$
        (\cref{met:g12:limcont:limits}).
  \item Geef de asymptoten van $f(x) = \dfrac{3x - 1}{x + 2}$
        (\cref{def:g12:limcont:asymptote}), met berekening van de betreffende
        limieten.
  \item Bereken $\lim_{x \to 3} \dfrac{x^2 - 9}{x - 3}$ (\omterm{def:g10:algebra:expand}{ontbinden}) en
        $\lim_{x \to 0} \dfrac{\sqrt{x + 1} - 1}{x}$ (toegevoegde
        uitdrukking).
  \item Met de insluitstelling (\cref{thm:g12:limcont:squeeze}): bereken
        $\lim_{x \to 0} x \sin\frac1x$ en
        $\lim_{x \to +\infty} \frac{\sin x}{x}$.
  \item Met samenstelling en overheersende termen:
        $\lim_{x \to +\infty} \dfrac{\sqrt{4x^2 + x}}{x}$.
\end{enumerate}

\medskip
\noindent\textbf{Deel II --- \omterm{def:g11:quad:discriminant}{Wortels}, met certificaat.}
\begin{enumerate}[resume]
  \item Toon aan dat de \omterm{def:g10:algebra:equation}{vergelijking} $x^3 + x = 5$ minstens één oplossing
        heeft in $\intcc{1}{2}$.
  \item Toon aan dat de oplossing uniek is op $\R$
        (\cref{thm:g12:limcont:bijection}).
  \item Bewijs het klassieke feit: \emph{elke veelterm van oneven graad heeft
        minstens één reële \omterm{def:g11:quad:discriminant}{wortel}}. (Schrijf het argument uit voor
        $p(x) = x^3 + a x^2 + b x + c$: bereken de twee oneindige limieten met
        de overheersende term en pas daarna de tussenwaardestelling toe op een
        groot \omterm{def:g10:numbers:interval}{interval}.)
  \item De \omterm{def:g11:quad:discriminant}{wortel} van vraag 6 met \omterm{thm:g12:limcont:ivt}{bisectie} lokaliseren: hoeveel halveringen
        van $\intcc{1}{2}$ waarborgen een fout onder $10^{-6}$
        (\cref{exo:g12:limcont:9})? Vergelijk met het verdubbelen van de
        decimalen bij de methode van Newton (\cref{pb:g11:deriv:1}): hoeveel
        stappen zou die ruwweg nodig hebben?
  \item Het tegenvoorbeeld dat de stelling bewaakt: $g(x) = \frac1x$ voldoet
        aan $g(-1) = -1 < 0$ en $g(1) = 1 > 0$, en verdwijnt toch nooit. Welke
        hypothese van de tussenwaardestelling faalt --- en wat leert het
        voorbeeld over het nagaan van hypothesen?
\end{enumerate}

\medskip
\noindent\textbf{Deel III --- Drie onmogelijk klinkende stellingen.}
\begin{enumerate}[resume]
  \item De stelling van het vaste punt op het lijnstuk: is
        $f : \intcc{0}{1} \to \intcc{0}{1}$ \omterm{def:g12:limcont:continuity}{continu}, bewijs dan dat een
        zekere $c$ voldoet aan $f(c) = c$. (Bestudeer $g(x) = f(x) - x$ in de
        twee uiteinden.)
  \item Vertaal vraag 11 naar de verfrommelde kaart: wat speelt de rol van
        $f$, waarom belandt ze in $\intcc{0}{1}$ (figuurlijk: binnen
        Frankrijk), en welk punt ligt precies boven zijn eigen plek? (Eén
        dimensie wordt stilzwijgend aangenomen --- de eerlijke
        tweedimensionale uitspraak is de stelling van Brouwer, universitaire
        volumes.)
  \item De evenaar: zij $T$ een continue temperatuur op een cirkel, opgevat
        als een $2\pi$-periodieke \omterm{def:g11:func:function}{functie} met $T(0) = T(2\pi)$. Stel
        $g(t) = T(t) - T(t + \pi)$. Bereken $g(0) + g(\pi)$, leid af dat $g$
        ergens op $\intcc{0}{\pi}$ verdwijnt, en formuleer de stelling over
        antipodale punten.
  \item De wiebelende tafel: een volkomen starre vierkante tafel staat op
        continue (zonder treden!) oneffen grond; drie poten raken altijd de
        grond, terwijl één zweeft op een hoogte $h(\theta)$ die \omterm{def:g12:limcont:continuity}{continu} van de
        draaihoek $\theta$ van de tafel afhangt. Leg uit waarom een
        kwartdraai de zwevende opening van teken doet veranderen --- en
        besluit dat een zekere hoek $\theta^*$ alle vier de poten stabiel
        zet. (Dit argument werd in 2005 als volwaardige stelling gepubliceerd;
        caféhouders kenden de truc al eerder.)
  \item De wandelaar van \cref{exo:g12:limcont:8} gebruikte hetzelfde schema.
        Formuleer in drie regels het sjabloon dat de vragen 11, 13 en 14 en de
        wandelaar delen --- welke \omterm{def:g11:func:function}{functie} bouw je, wat ga je in de uiteinden
        na, en wat roep je in?
  \item Wat de tussenwaardestelling \emph{niet} geeft: het bestaan kwam
        gratis, maar welke twee vragen laat ze open, en welke gereedschappen
        uit deze opgave beantwoorden ze (vragen 7 en 9)?
\end{enumerate}

\medskip
\noindent\textbf{Deel IV --- Portretten met asymptoten.}
\begin{enumerate}[resume]
  \item Herschrijf $f(x) = \dfrac{3x - 1}{x + 2}$ als
        $3 - \dfrac{7}{x + 2}$ en beschrijf daarmee de kromme volledig:
        asymptoten, verloop aan elke kant, schets. (Een verschoven \omterm{def:g11:func:reference}{hyperbool}
        --- de kromme van de \omterm{def:g11:stat:mean}{gemiddelde} kostprijs uit
        \cref{pb:g10:reffunc:1} was er ook een.)
  \item De \omterm{def:g11:func:function}{functie} $f(x) = \dfrac{x^2 - 1}{x - 1}$ is in $1$ niet
        gedefinieerd. Welke waarde voor $f(1)$ maakt haar daar \omterm{def:g12:limcont:continuity}{continu} --- en
        waarom noemen analisten zo'n punt een \emph{ophefbare}
        discontinuïteit?
  \item De entierfunctie $x \mapsto \lfloor x \rfloor$ (het grootste gehele
        getal $\leq x$): waar is ze \omterm{def:g12:limcont:continuity}{continu}, waar springt ze, en waarom faalt
        het besluit van de tussenwaardestelling voor haar (zoek $a < b$ met
        $\lfloor a \rfloor < \frac12 < \lfloor b \rfloor$ terwijl
        $\lfloor x \rfloor = \frac12$ geen oplossing heeft)?
  \item Slotstuk --- de twee gezichten van de \omterm{def:g12:limcont:continuity}{continuïteit}: de plaatselijke
        belofte (waarde $=$ limiet: geen teleportatie) en de globale kracht
        (de tussenwaardestelling: elke tussenliggende waarde wordt bezocht).
        Breng de vier wonderen uit deel III onder bij het juiste gezicht, en
        geef de leuze over het oversteken van de rivier haar precieze
        wiskundige naam.
\end{enumerate}
\end{problem}
